% Appendix A for Final Report
% Software Quality Metrics Analysis of Open Source Projects

\appendix

\section{Phase 1 Planned Configuration (GitHub Actions)}

This section documents the original Phase 1 approach that involved
integrating GitHub Actions workflows into target repositories for
automated static code analysis and quality metric collection.

\subsection{Planned GitHub Actions Workflow}

The intended workflow would have been added to target repositories at
\texttt{.github/workflows/quality-metrics.yml}:

\begin{lstlisting}[language=yaml, caption=Planned quality-metrics.yml Workflow]
name: Software Quality Metrics Collection

on:
  push:
    branches: [main, develop]
  pull_request:
    branches: [main]

jobs:
  quality-analysis:
    runs-on: ubuntu-latest

    steps:
      - name: Checkout code
        uses: actions/checkout@v3

      - name: Set up Python
        uses: actions/setup-python@v4
        with:
          python-version: '3.11'

      - name: Install dependencies
        run: |
          pip install radon pytest pytest-cov ruff mypy
          pip install -r requirements.txt

      - name: Run Radon - LOC Analysis
        run: |
          radon raw . -j > metrics/radon_loc.json
          radon cc . -j > metrics/radon_complexity.json
          radon mi . -j > metrics/radon_maintainability.json

      - name: Run Pytest with Coverage
        run: |
          pytest --cov=. --cov-report=json \
                 --cov-report=term \
                 --junitxml=metrics/pytest_results.xml

      - name: Run Ruff Linter
        run: |
          ruff check . --output-format=json \
               > metrics/ruff_violations.json

      - name: Run Mypy Type Checker
        run: |
          mypy . --json-report metrics/mypy_report

      - name: Enforce Quality Gates
        run: |
          python .github/scripts/enforce_quality_gates.py

      - name: Publish Metrics
        uses: actions/upload-artifact@v3
        with:
          name: quality-metrics
          path: metrics/

      - name: Send to Observability Platform
        run: |
          python .github/scripts/publish_to_otel.py
        env:
          OTEL_ENDPOINT: ${{ secrets.OTEL_ENDPOINT }}
          OTEL_API_KEY: ${{ secrets.OTEL_API_KEY }}
\end{lstlisting}

\subsection{Planned Quality Gate Enforcement}

The workflow would have included automated quality gates to prevent
merging code that failed thresholds:

\begin{lstlisting}[language=Python, caption=enforce\_quality\_gates.py]
import json
import sys

# Load metrics
with open('metrics/radon_complexity.json') as f:
    complexity = json.load(f)

with open('metrics/pytest_coverage.json') as f:
    coverage = json.load(f)

with open('metrics/ruff_violations.json') as f:
    linting = json.load(f)

# Define thresholds
THRESHOLDS = {
    'max_complexity': 15,
    'min_coverage': 70,
    'max_linting_violations': 50,
    'min_maintainability': 50
}

# Check thresholds
violations = []

# Check cyclomatic complexity
for module in complexity:
    for func in module.get('functions', []):
        if func['complexity'] > THRESHOLDS['max_complexity']:
            violations.append(
                f"Function {func['name']} has complexity "
                f"{func['complexity']} (max: {THRESHOLDS['max_complexity']})"
            )

# Check test coverage
total_coverage = coverage['totals']['percent_covered']
if total_coverage < THRESHOLDS['min_coverage']:
    violations.append(
        f"Test coverage {total_coverage}% is below "
        f"minimum {THRESHOLDS['min_coverage']}%"
    )

# Check linting violations
violation_count = len(linting)
if violation_count > THRESHOLDS['max_linting_violations']:
    violations.append(
        f"Linting violations {violation_count} exceed "
        f"maximum {THRESHOLDS['max_linting_violations']}"
    )

# Report violations
if violations:
    print("Quality gate FAILED:")
    for v in violations:
        print(f"  - {v}")
    sys.exit(1)
else:
    print("Quality gate PASSED")
    sys.exit(0)
\end{lstlisting}

\subsection{Planned Tool Configurations}

\subsubsection{Radon Configuration}

Radon would analyze code complexity without additional configuration,
but the following commands were planned:

\begin{lstlisting}[language=bash, caption=Radon Analysis Commands]
# Lines of Code analysis
radon raw . --json --output-file=radon_loc.json

# Cyclomatic Complexity
radon cc . --json --min=C --output-file=radon_complexity.json

# Maintainability Index
radon mi . --json --min=B --output-file=radon_maintainability.json

# Halstead metrics (code effort estimation)
radon hal . --json --output-file=radon_halstead.json
\end{lstlisting}

\subsubsection{Pytest Configuration}

\begin{lstlisting}[language=ini, caption=pytest.ini]
[pytest]
testpaths = tests
python_files = test_*.py
python_classes = Test*
python_functions = test_*

# Coverage settings
addopts =
    --cov=src
    --cov-report=html
    --cov-report=json
    --cov-report=term-missing
    --cov-fail-under=70
    --junitxml=junit.xml
    --verbose

# Coverage exclusions
[coverage:run]
omit =
    */tests/*
    */venv/*
    */migrations/*
    */__pycache__/*

[coverage:report]
precision = 2
show_missing = True
skip_covered = False
\end{lstlisting}

\subsubsection{Ruff Configuration}

\begin{lstlisting}[language=toml, caption=pyproject.toml - Ruff Configuration]
[tool.ruff]
line-length = 100
target-version = "py311"

select = [
    "E",   # pycodestyle errors
    "W",   # pycodestyle warnings
    "F",   # pyflakes
    "I",   # isort
    "C",   # flake8-comprehensions
    "B",   # flake8-bugbear
    "UP",  # pyupgrade
    "N",   # pep8-naming
]

ignore = [
    "E501",  # line too long (handled by formatter)
    "B008",  # function call in argument defaults
]

[tool.ruff.per-file-ignores]
"__init__.py" = ["F401"]  # unused imports
"tests/*.py" = ["S101"]   # allow assert statements

[tool.ruff.isort]
known-first-party = ["src"]
\end{lstlisting}

\subsubsection{Mypy Configuration}

\begin{lstlisting}[language=ini, caption=mypy.ini]
[mypy]
python_version = 3.11
warn_return_any = True
warn_unused_configs = True
disallow_untyped_defs = True
disallow_incomplete_defs = True
check_untyped_defs = True
no_implicit_optional = True
warn_redundant_casts = True
warn_unused_ignores = True
warn_no_return = True
strict_equality = True
show_error_codes = True

# Per-module options
[mypy-tests.*]
disallow_untyped_defs = False

[mypy-third_party_lib.*]
ignore_missing_imports = True
\end{lstlisting}

\subsection{Planned Semantic Versioning Integration}

The workflow would have used \texttt{python-semantic-release} for
automated version management:

\begin{lstlisting}[language=yaml, caption=Semantic Release Workflow]
name: Semantic Release

on:
  push:
    branches: [main]

jobs:
  release:
    runs-on: ubuntu-latest
    steps:
      - uses: actions/checkout@v3
        with:
          fetch-depth: 0

      - name: Python Semantic Release
        uses: relekang/python-semantic-release@master
        with:
          github_token: ${{ secrets.GITHUB_TOKEN }}
          pypi_token: ${{ secrets.PYPI_TOKEN }}
\end{lstlisting}

\begin{lstlisting}[language=toml, caption=pyproject.toml - Semantic Release Config]
[tool.semantic_release]
version_variable = "src/__init__.py:__version__"
version_source = "tag"
commit_version_number = true
tag_commit = true
upload_to_pypi = false
upload_to_release = true

branch = "main"
hvcs = "github"

# Commit parsing
commit_parser = "angular"
major_on_zero = true
allow_zero_version = true

# Version bumping rules
[tool.semantic_release.commit_parser_options]
allowed_tags = [
    "feat",      # New feature (minor bump)
    "fix",       # Bug fix (patch bump)
    "docs",      # Documentation only
    "style",     # Code style changes
    "refactor",  # Code refactoring
    "perf",      # Performance improvements
    "test",      # Test additions
    "build",     # Build system changes
    "ci",        # CI configuration
]

major_tags = ["BREAKING CHANGE", "BREAKING"]
minor_tags = ["feat"]
patch_tags = ["fix", "perf"]
\end{lstlisting}

\section{Data Collection Scripts}

The data collection process involved manually extracting data from
GitHub's web interface and organizing it into CSV format. While no
automated scraping scripts were used due to permission constraints,
the following describes the manual collection workflow:

\subsection{GitHub Web UI Navigation}

For each repository (LangChain and Intel XPU Backend), data was
collected through:

\begin{enumerate}
  \item \textbf{Pull Request Data Collection}
    \begin{itemize}
      \item Navigate to repository $\rightarrow$ Pull Requests tab
      \item Filter by date range (e.g., October 29 - November 27,
        2025 for LangChain)
      \item Filter by state: \texttt{is:pr is:open}, \texttt{is:pr
        is:merged}, \texttt{is:pr is:closed}
      \item Copy PR titles, numbers, and timestamps to CSV
    \end{itemize}

  \item \textbf{Issue Data Collection}
    \begin{itemize}
      \item Navigate to repository $\rightarrow$ Issues tab
      \item Filter by creation date within analysis window
      \item Identify bug-related issues using keyword search: bug,
        error, fix, broken, crash, fail, regression, exception
      \item Export issue titles and creation timestamps
    \end{itemize}

  \item \textbf{Repository Metadata}
    \begin{itemize}
      \item Navigate to Insights $\rightarrow$ Pulse for activity summary
      \item Record commit hashes and dates from commit history
      \item Note relative timestamps ("2 days ago") and convert to
        absolute dates
    \end{itemize}
\end{enumerate}

\subsection{CSV Data Organization}

Data was organized into four CSV files per repository:

\begin{lstlisting}[language=bash, caption=CSV File Structure]
# opened_requests.csv
msg,hash,date_rel,date_abs
"Add authentication feature",#12345,"2 days ago","2025-11-27"

# merged_requests.csv
msg,hash,date_rel,date_abs
"Fix authentication bug",#12346,"1 day ago","2025-11-28"

# closed_requests.csv
msg,hash,date_rel,date_abs
"Update documentation",#12347,"3 hours ago","2025-11-29"

# opened_issues.csv
msg,hash,date_rel,date_abs
"Bug: Login fails on Safari",#678,"1 week ago","2025-11-22"
\end{lstlisting}

\section{Analysis Source Code}

\subsection{Key Functions}

\begin{itemize}
  \item \texttt{load\_data(filepath)}: Loads CSV data and
    preprocesses timestamps
    \begin{itemize}
      \item Parses date\_abs column into datetime objects
      \item Handles missing or malformed timestamps
      \item Returns pandas DataFrame
    \end{itemize}

  \item \texttt{calculate\_deployment\_frequency(merged\_df, days)}:
    Calculates DF metric
    \begin{itemize}
      \item Counts merged PRs per day over analysis window
      \item Computes mean, median, and weekly aggregate
      \item Classifies performance: Elite ($>$1/day), High (weekly),
        Medium (monthly), Low ($<$monthly)
    \end{itemize}

  \item \texttt{calculate\_lead\_time(opened\_df, merged\_df)}:
    Calculates LT metric
    \begin{itemize}
      \item Matches opened PRs with merged PRs using temporal proximity
      \item Computes time delta in days
      \item Returns median, mean, 25th/75th/90th percentiles
      \item Classifies: Elite ($<$1 day), High (1-7 days), Medium
        (1-4 weeks), Low ($>$1 month)
    \end{itemize}

  \item \texttt{calculate\_change\_failure\_rate(issues\_df,
    merged\_df)}: Calculates CFR metric
    \begin{itemize}
      \item Identifies bug issues using keyword matching on issue
        titles/descriptions
      \item Counts bug issues within analysis window
      \item Computes ratio: (bug issues / total deployments) $\times$ 100\%
      \item Classifies: Elite (0-15\%), High (16-30\%), Medium
        (31-45\%), Low ($>$45\%)
    \end{itemize}

  \item \texttt{calculate\_time\_to\_restore(issues\_df,
    merged\_df)}: Calculates MTTR metric
    \begin{itemize}
      \item Attempts to pair bug issues with fix PRs
      \item Searches for issue references in PR titles/bodies
      \item Computes time from bug opened to fix merged
      \item Returns median restoration time
      \item Classifies: Elite ($<$1 hour), High ($<$1 day), Medium
        (1-7 days), Low ($>$1 week)
    \end{itemize}

  \item \texttt{generate\_visualizations(metrics\_dict,
    output\_dir)}: Creates charts
    \begin{itemize}
      \item Deployment frequency time series
      \item Lead time distribution histogram
      \item Change failure rate comparison (deployments vs bugs)
      \item Performance classification summary
    \end{itemize}
\end{itemize}

\subsection{Usage Example}

\begin{lstlisting}[language=Python, caption=Running DORA Metrics Analysis]
import pandas as pd
from dora_metrics import *

# Load data
opened = load_data('data/langchain/opened_requests.csv')
merged = load_data('data/langchain/merged_requests.csv')
closed = load_data('data/langchain/closed_requests.csv')
issues = load_data('data/langchain/opened_issues.csv')

# Calculate metrics
df = calculate_deployment_frequency(merged, days=30)
lt = calculate_lead_time(opened, merged)
cfr = calculate_change_failure_rate(issues, merged)
mttr = calculate_time_to_restore(issues, merged)

# Generate visualizations
metrics = {'df': df, 'lt': lt, 'cfr': cfr, 'mttr': mttr}
generate_visualizations(metrics, 'output/langchain/')

# Print summary
print(f"Deployment Frequency: {df['mean']:.2f}/day")
print(f"Lead Time (median): {lt['median']:.2f} days")
print(f"Change Failure Rate: {cfr['rate']:.2f}%")
print(f"Time to Restore: {mttr['median']:.2f} hours")
\end{lstlisting}

\section{Raw Data Samples}

\subsection{LangChain Repository Data}

Sample entries from the LangChain analysis dataset (October 29 -
November 27, 2025):

\begin{table}[H]
  \centering
  \caption{Sample Merged Pull Requests (LangChain)}
  \label{tab:langchain-sample-prs}
  \small
  \begin{tabular}{@{}llll@{}}
    \toprule
    \textbf{PR \#} & \textbf{Title} & \textbf{Date Relative} &
    \textbf{Date Absolute} \\ \midrule
    \#25843 & anthropic[patch]: Release 0.41.1 & 2 days ago & 2025-11-25 \\
    \#25841 & docs: Update Anthropic tools example & 2 days ago & 2025-11-25 \\
    \#25839 & core[patch]: Fix streaming in RunnableBinding & 3 days
    ago & 2025-11-24 \\
    \#25832 & community[patch]: Add timeout to Redis & 4 days ago &
    2025-11-23 \\
    \#25829 & docs: Add citation format to papers & 5 days ago &
    2025-11-22 \\ \bottomrule
  \end{tabular}
\end{table}

\begin{table}[H]
  \centering
  \caption{Sample Bug Issues (LangChain)}
  \label{tab:langchain-sample-issues}
  \small
  \begin{tabular}{@{}llll@{}}
    \toprule
    \textbf{Issue \#} & \textbf{Title} & \textbf{Date Relative} &
    \textbf{Date Absolute} \\ \midrule
    \#4521 & Bug: Streaming fails with Azure OpenAI & 1 day ago & 2025-11-26 \\
    \#4518 & Error when using callback with agents & 3 days ago & 2025-11-24 \\
    \#4510 & Fix: Memory leak in conversation buffer & 5 days ago &
    2025-11-22 \\
    \#4502 & Regression: Tool calling broken in 0.41.0 & 1 week ago &
    2025-11-20 \\ \bottomrule
  \end{tabular}
\end{table}

\subsection{Intel XPU Backend Data}

Sample entries from Intel XPU Backend analysis dataset (November 4 -
December 2, 2025):

\begin{table}[H]
  \centering
  \caption{Sample Merged Pull Requests (Intel XPU Backend)}
  \label{tab:intel-sample-prs}
  \small
  \begin{tabular}{@{}llll@{}}
    \toprule
    \textbf{PR \#} & \textbf{Title} & \textbf{Date Relative} &
    \textbf{Date Absolute} \\ \midrule
    \#2104 & Support for Triton 3.2.0 & 1 day ago & 2025-12-01 \\
    \#2101 & Fix GEMM kernel for FP16 & 2 days ago & 2025-11-30 \\
    \#2098 & Add support for grouped convolutions & 4 days ago & 2025-11-28 \\
    \#2095 & Optimize reduction kernels & 6 days ago & 2025-11-26 \\
    \#2092 & Fix memory allocation in compiler & 1 week ago &
    2025-11-25 \\ \bottomrule
  \end{tabular}
\end{table}

\section{OpenTelemetry Configuration Files}

\subsection{Docker Compose Configuration}

The Phase 2 OpenTelemetry approach used the following Docker Compose setup:

\begin{lstlisting}[language=yaml, caption=docker-compose.yml for OTel Stack]
version: '3.8'

services:
  otel-collector:
    image: otel/opentelemetry-collector-contrib:latest
    container_name: otel-collector
    command: ["--config=/etc/otel-collector-config.yaml"]
    volumes:
      - ./otel-collector-config.yaml:/etc/otel-collector-config.yaml
    ports:
      - "4317:4317"   # OTLP gRPC receiver
      - "4318:4318"   # OTLP HTTP receiver
      - "8888:8888"   # Prometheus metrics
    networks:
      - otel-network
    depends_on:
      - openobserve

  openobserve:
    image: public.ecr.aws/zinclabs/openobserve:latest
    container_name: openobserve
    ports:
      - "5080:5080"
    environment:
      - ZO_ROOT_USER_EMAIL=admin@example.com
      - ZO_ROOT_USER_PASSWORD=admin123
      - ZO_DATA_DIR=/data
    volumes:
      - openobserve-data:/data
    networks:
      - otel-network

networks:
  otel-network:
    driver: bridge

volumes:
  openobserve-data:
\end{lstlisting}

\subsection{OpenTelemetry Collector Configuration}

\begin{lstlisting}[language=yaml, caption=otel-collector-config.yaml]
receivers:
  github:
    github_org: langchain-ai
    search_query: "org:langchain-ai is:pr"
    scrapers:
      github:
        collection_interval: 60s
    auth:
      authenticator: bearertokenauth

extensions:
  bearertokenauth:
    token: "${GITHUB_TOKEN}"

processors:
  batch:
    timeout: 10s
    send_batch_size: 100

exporters:
  otlphttp:
    endpoint: "http://openobserve:5080/api/default/"
    headers:
      Authorization: "Basic ${OPENOBSERVE_CREDENTIALS}"
      stream-name: "github-metrics"

service:
  extensions: [bearertokenauth]
  pipelines:
    metrics:
      receivers: [github]
      processors: [batch]
      exporters: [otlphttp]
\end{lstlisting}

\subsection{Performance Classification Logic}

\begin{lstlisting}[language=Python, caption=DORA Classification]
def classify_deployment_frequency(deployments_per_day):
    """Elite: >1/day, High: weekly, Medium: monthly, Low: <monthly"""
    if deployments_per_day >= 1:
        return 'Elite'
    elif deployments_per_day >= 1/7:
        return 'High'
    elif deployments_per_day >= 1/30:
        return 'Medium'
    else:
        return 'Low'

def classify_lead_time(median_days):
    """Elite: <1 day, High: 1-7 days, Medium: 1-4 weeks, Low: >1 month"""
    if median_days < 1:
        return 'Elite'
    elif median_days <= 7:
        return 'High'
    elif median_days <= 28:
        return 'Medium'
    else:
        return 'Low'

def classify_change_failure_rate(percentage):
    """Elite: 0-15%, High: 16-30%, Medium: 31-45%, Low: >45%"""
    if percentage <= 15:
        return 'Elite'
    elif percentage <= 30:
        return 'High'
    elif percentage <= 45:
        return 'Medium'
    else:
        return 'Low'

def classify_time_to_restore(median_hours):
    """Elite: <1 hour, High: <1 day, Medium: 1-7 days, Low: >1 week"""
    if median_hours < 1:
        return 'Elite'
    elif median_hours < 24:
        return 'High'
    elif median_hours < 168:  # 7 days
        return 'Medium'
    else:
        return 'Low'
\end{lstlisting}

\end{document}
